\section{History} 
\label{sec:history}

The history of \acp{bss} reflects over half a century of experimentation, learning, and innovation in response to the challenges of providing shared bicycle access in urban environments (Figure~\ref{fig:table_generations}).

\begin{figure}[htp!]
    \centering
    \includegraphics[width=0.8\textwidth]{0_plots/figure_generations.pdf}
    \caption{\ac{bss} generations and their main characteristics. Innovations are shown in black, and features inherited from previous generations are shown in grey. Author’s own elaboration.}
    \label{fig:table_generations}
\end{figure}

The earliest recorded initiative was the \textit{White Bike Plan}~\cite{Shirky2008, vanderZee2016}, launched in Amsterdam in 1965 by the Provos, a countercultural group advocating for environmental and anti-car urban policies. The initiative involved painting 50 bicycles white and distributing them, unlocked and free of charge, throughout the city center. Despite its symbolic value, the initiative failed shortly after launch due to theft, vandalism, and police intervention. Similar free bike systems appeared in the following decades in cities such as La Rochelle, France in 1976 with the \textit{Vélos Jaunes} program~\cite{Parbaud2023}, and in Cambridge, U.K. in 1993 with the \textit{Green Bike Scheme}~\cite{Pilgrim2017}. These efforts, often operated by municipalities or non-profits, prioritized accessibility and environmental benefits but offered little in terms of system control, user accountability, or reliability~\cite{Medard2019}. This era is now referred to as the \textbf{first generation} of bike-sharing.

Recognizing the limitations of unregulated access, the \textbf{second generation} introduced coin-deposit mechanisms and designated docking infrastructure. In 1995, Copenhagen launched \textit{Bycykler København}~\cite{Walker2018}, a large-scale program with 1,100 specially designed bicycles that could be unlocked with a coin and returned to designated stations. This approach reduced theft and introduced a modest form of user accountability. Similar programs followed in cities of Europe such as \textcolor{red}{Sandnes, Norway in 1996, Helsinki, Finland in 2000, and Aarhus, Denmark in 2005}. In the US, multiple cities also implemented similar programs such as the \textit{Yellow Bike Project} in Minneapolis and St. Paul in 1998~\cite{Shaheen2010} and the \textit{Red Bikes} program in Madison in 2000~\cite{Charly2012}. These systems still operated without digital technologies, and while they were more structured than their predecessors, they lacked mechanisms to enforce time limits or ensure returns. As a result, bikes were frequently used for extended periods or disappeared altogether, making system maintenance costly and unsustainable~\cite{Shaheen2010, Medard2019}.

A major shift in bike-sharing occurred with the emergence of the \textbf{third generation}. These systems were generally defined by four key innovations~\cite{Shaheen2010}: (i) visually distinctive bicycles build on purpose to discourage theft and simplify maintenance; (ii) designated docking infrastructure where users could pick up and return bikes; (iii) on-site user interfaces or kiosks to handle rentals and returns; and (iv) access systems based on magnetic cards, smartcards, or other digital credentials. Thanks to the introduced digital infrastructure, \acp{bss} managers were now able to ensure a more secure system by monitoring usage and user information. It is important to note that, from the third generation onward, large-scale deployments often relied on public–private partnerships with advertising or technology firms (e.g., JCDecaux in Paris and Clear Channel in Barcelona’s first \textit{Bicing} contract), which played a central role in financing and operations.

Although second-generation systems had already been deployed outside Europe, the third generation marked the first truly global wave of bike-sharing on the 2000s onward, with systems appearing across Europe, North and South America, and Asia~\cite{Shaheen2010}. One early third generation example was the \textit{Vélo à la Carte} system launched in Rennes in 1998, which introduced smartcard-based access and automated docking stations~\cite{DeMaio2009}. However, it was the roll-out of \textit{Vélo’v} in Lyon in 2005, and specially the launch of \textit{Vélib’} in Paris in 2007 that really popularized this new generation of bike-sharing~\cite{DeMaio2009}. Vélib’ became a global reference point with over 20,000 bicycles and 1,400 stations across the city, operating 24/7 with a subscription model and a progressive pricing structure that incentivized short trips~\cite{Shaheen2010}. These systems demonstrated the scalability and attractiveness of bike-sharing when supported by robust infrastructure and public–private partnerships.

Beyond Europe, third-generation \acp{bss} began to gain momentum in North America, South America, and Asia during the late 2000s. In the United States, Washington D.C.'s \textit{SmartBike}—launched in 2008—was among the first IT-based programs~\cite{Silverman2008}, though its limited scale and visibility were soon eclipsed by more ambitious efforts. That same year, Montreal introduced \textit{BIXI}~\cite{Entreprisesàvendre2025, Shaheen2010}, a modular and open-data-based system that served as a technological model for other cities and was later adopted in New York City (\textit{Citi Bike})~\cite{RudinCenter2015}, London (\textit{Santander Cycles})~\cite{BBC2020}, and Boston (\textit{Bluebikes})~\cite{BluebikesAbout}. In South America there were a few examples of third-generation systems. For example, Rio de Janeiro's \textit{Samba}, Sao Paulo's \textit{UseBike}~\cite{Shaheen2010}, and Santiago de Chile's \textit{Bikesantiago}~\cite{Shaheen2010}. Meanwhile, Asian cities quickly embraced large-scale systems. Hangzhou’s Public Bicycle program, launched in 2008~\cite{Guzman2011}, grew to become one of the largest in the world, with tens of thousands of bikes and an extensive terminal network. Notably, at this point in time many of these cities skipped earlier generations entirely, opting from the outset for smart, digitally managed systems that could be integrated into broader urban mobility strategies.

The concept of a \textbf{fourth generation} of \acp{bss} was first introduced by DeMaio~\cite{DeMaio2009} on 2009, who envisioned future systems as more efficient, sustainable, and user-friendly through innovations such as improved bicycle rebalancing, mobile or off-grid station infrastructure, \acp{e-b}, and novel business models. Building on this vision, Shaheen et al.~\cite{Shaheen2010} formalized the idea as an emerging paradigm extending third-generation systems. They identified four defining features: (i) flexible or self-powered docking infrastructure (e.g., mobile or solar-powered stations); (ii) advanced bicycle redistribution strategies, including automated rebalancing and user-based incentives for easing rebalancing; (iii) integration with other transport modes via smartcards; (iv) enhanced technology such as GPS tracking, touchscreen kiosks, and the incorporation of \acp{e-b}. Among these, the diffusion of \acp{e-b} and the increasing integration of digital platforms and data-driven management have proven particularly transformative, reducing physical effort and enabling longer trips, allowing real-time monitoring, and broadening the appeal of bike-sharing to a wider range of users.

While Shaheen et al.~\cite{Shaheen2010} originally defined these criteria in 2010, subsequent studies have tended to focus selectively on certain subsets of these traits. For instance, Zheng et al.~\cite{Zheng2020} provide the most comprehensive formulation, retaining all four original elements. Fishman~\cite{Fishman2016}, by contrast, places particular emphasis on dockless operation, the rise of \acp{e-b}, and smartcard integration, while Chen et al.~\cite{Chen2020} similarly concentrate on app-based dockless systems. Meanwhile, Meireles and Silva~\cite{Meireles2013} stress the integration of \acp{e-b} and system efficiency from both user and operator perspectives. Taken together, these contributions illustrate how the notion of a fourth generation has evolved into a flexible conceptual framework rather than a strict category. In contrast to earlier generations, which were defined by one or two core innovations, the fourth generation brings together multiple components, making full implementation relatively rare in practice. Prominent examples of fourth-generation systems include Paris' 2018 \textit{Vélib’ Métropole} version, with \acp{e-b} and temporary solar-powered stations~\cite{Garrigues2017}, and Barcelona’s \textit{Bicing}, which incorporated a large fleet of \acp{e-b} and an app-based digital access~\cite{Bicing2025}.

This ambiguity also explains why some scholars propose a \textbf{fifth generation}—viewed as an extension of the fourth generation’s emphasis on digital integration—, in which dockless operation itself becomes the defining feature~\cite{Xu2024, Vallez2021, Guidon2019, Chen2018}. These schemes eliminate the need for fixed infrastructure, relying instead on GPS-enabled bicycles and smartphone apps for locating, unlocking, and paying for rides. First introduced at scale by \textit{Ofo} in China~\cite{Zhang2018} and soon joined by competitors such as \textit{Mobike}, this model was later expanded internationally by companies including \textit{Lime}, \textit{Call-a-bike}, and others, offering new conveniences: users no longer need to secure a vacant dock at their destination, and trips can start and end with greater spatial flexibility. Yet this flexibility introduces new challenges. The absence of docking infrastructure has led to problems of sidewalk clutter, parking regulation, and public space encroachment—issues especially visible in the early large-scale Chinese deployments. In cities such as Shanghai and Beijing, overflowing bikes blocked pedestrian walkways and spilled into residential areas, prompting public backlash and municipal regulation~\cite{Wang2019}.

\begin{figure}[htp!]
    \centering
    \includegraphics[width=0.8\textwidth]{0_plots/bss_growth_worldwide/bss_growth_worldwide.png}
    \caption{Growth of \acp{bss} worldwide. Data sourced from the Meddin Bike-sharing World Map Report of 2022~\cite{MeddinWorldMapReport2022}. Author’s own elaboration.}
    \label{fig:bss_growth_worldwide}
\end{figure}

As shown in Figure~\ref{fig:bss_growth_worldwide}, \acp{bss} have grown from a handful of initiatives in the 1990s to more than 2,000 operating systems worldwide by 2022. Moreover, although the COVID-19 pandemic caused a temporary dip in usage in 2020, many systems rebounded quickly as cycling was perceived as a safe and resilient mode of transport during health crises \cite{Teixeira2024}. By 2025, there are 2,100 active systems worldwide, with nearly 40\% of these systems including \acp{e-b}, and the global fleet exceeds 10 million shared bikes~\cite{MeddinWorldMap2025}. This rapid expansion reflects not only technological progress, but also the growing recognition of shared micromobility as a viable component of sustainable urban transport. 